% !PRINTED 644
\weberpage{644}
%<<CONT
zum Theil Null sein können. Durch dieselbe Veränderung ergiebt
 dann \S.\,185 (6) den Werth von $\eta_0$.

Damit sind also die beschränkenden Voraussetzungen des
\S.\,185 beseitigt, und wir haben das allgemeine Theorem:

\gesperrt{I. Jede Wurzel $\xi$ einer metacyklischen Gleichung
vom Primzahlgrad $n$ kann in der Form dargestellt
werden}
\weq{9}{\xi = A + \sum_{0,\,n-2}^{\nu} K_\nu \tau_\nu^{r_{n-2}} \tau_{\nu+1}^{r_{n-3}} \cdots \tau_{\nu+n-2}^{r_0},}
\gesperrt{worin $A$ eine rationale Grösse, $k_0, k_1 \dots k_{n-2}$ die
von einander und von Null verschiedenen Wurzeln einer cyklischen
 Gleichung $(n-1)^{\text{ten}}$ Grades, $K_\nu$ eine rationale Function
von $k_\nu$ ist, deren Form für alle $\nu$ dieselbe ist. Die Exponenten
$r_0, r_1 \dots r_{n-2}$ sind die kleinsten positiven Reste der
Zahlen $1, g, g^2 \dots g^{n-2}$, wenn $g$ eine primitive
Wurzel von $n$ ist. Die $n$ Werthe, die man aus (9)
erhält, wenn man den Radicalen $\tau_\nu = \sqrt[n]{k_\nu}$ ihre verschiedenen
 Werthe beilegt, sind die $n$ Wurzeln
einer und derselben rationalen Gleichung.}

Die Formel (9) ergiebt sich aus \S.\,185 (6), wenn $A$ und $K_\nu$
durch $nA$ und $nK_\nu$ ersetzt wird, was zur Vereinfachung geschehen
 ist.

Das Theorem I. lässt sich nun aber auch umkehren.

Um das nachzuweisen, bezeichnen wir mit $\varepsilon$ irgend eine
imaginäre $n^{\text{te}}$ Einheitswurzel und setzen
\weq{10}{\begin{aligned}\xi_h &= A + \sum_{0,\,n-2}^{\nu} \varepsilon^{h r_\nu} K_\nu \tau_\nu^{r_{n-2}} \tau_{\nu+1}^{r_{n-3}} \cdots \tau_{\nu+n-2}^{r_0}\\ h &= 0, 1, 2 \dots n-1,\end{aligned}}
oder abgekürzt
\weq{11}{\xi_h = A + \sum_{0,\,n-2}^{\nu} \varepsilon^{h r_\nu} K_\nu \sqrt[n]{R_\nu},}
worin wie früher
\weq{12}{\sqrt[n]{R_\nu} = \tau_\nu^{r_{n-2}} \tau_{\nu+1}^{r_{n-3}} \cdots \tau_{\nu+n-2}^{r_0}}
gesetzt ist. Wir haben nun die Aenderungen zu untersuchen,
die sich für $\xi_h$ ergeben, wenn wir

1. eines der Radicale $\tau_\nu = \sqrt[n]{k_\nu}$ anders bestimmen und

2. die $k_0, k_1 \dots k_{n-2}$ cyklisch vertauschen.
